% Event rules; encoding and setup details are in Appendix A.
공통 검사에 더할 회계 규칙은 거래 종류로 정한다. Entry·Transfer·
Proceed·Recall·Merge·Split·Exit에는 종류별 고정 검증기를 쓴다.
Process와 Issue는 정책 식별자 $p$로 등록된 검증기와 입출력 형식을
선택한다. Process에는 해당 정책을 사용할 권한도 필요하다.
각 거래는 최종 증명 하나를 제출한다.

원장은 증명 외에도 입력 등록·소유권·미소비·활성 상태, 중복 입력,
새 출력, 정책·권한·기한을 검사한다. 모두 통과해야 입력 소비,
출력 생성과 감사 기록을 한 거래에서 함께 갱신한다.
정확한 공개 입력과 설정·수락 절차는 부록~\ref{app:acceptance}에 둔다.

Table~\ref{tab:events}는 아홉 거래의 입출력을 보여 준다.
$N,V,C$는 Note·Voucher·Claim, $N_C$는 송신자에게 남는 Note다.
아래 식은 이 논문에서 선택한 참조 회계 규칙이다. 다른 등록 정책에는
그 정책이 실제로 검사한 조건만 적용된다.
\begin{table}[t]
\centering\small
\caption{Ordered event inputs and outputs}
\label{tab:events}
\begin{tabularx}{\columnwidth}{@{}lYrr@{}}
\toprule
Event & Transition & Parents & Future IDs \\
\midrule
Entry & $\varnothing\to N$ & 0 & 1 \\
Transfer & $N\to(V,N_C)$ & 1 & 2 \\
Proceed & $V\to N_R$ & 1 & 1 \\
Recall & $V\to N_S$ & 1 & 1 \\
Merge & $(N_1,N_2)\to N$ & 2 & 1 \\
Split & $N\to(N_1,N_2)$ & 1 & 2 \\
Process & $N^m\to N^n$ & $m$ & $n$ \\
Exit & $N\to\varnothing$ & 1 & 0 \\
Issue & $N\to C$ & 1 & 0 \\
\bottomrule
\end{tabularx}
\end{table}

\subsection{등록과 이전}

초기 등록(Entry)은 외부에서 인정한 근거로 일반용 Note를 만든다.
입력은 없고 $q>0$, $0\leq a\leq q$, $e=e_0\geq0$이어야 하며,
등록 승인자의 승인이 필요하다. 생성자가 최초 소유자다.
초기 탄소 생략 시의 값은 부록~\ref{app:acceptance}에 정한다.

이전(Transfer)은 일반용 Note를 소비해 보내는 Voucher와
남기는 Note를 만든다. 두 출력의 문서·역할은 입력과 같고,
Voucher에는 수신자가 고정된다.
아래첨자 $\mathrm{in},T,C$는 입력·보내는 몫·남기는 몫이다.
\begin{align}
 q_{\mathrm{in}}&=q_T+q_C,\quad q_T>0,\quad q_C\geq0,\nonumber\\
 a_C&=\left\lfloor\frac{a_{\mathrm{in}}q_C}{q_{\mathrm{in}}}\right\rfloor,
 \quad a_T=a_{\mathrm{in}}-a_C,\nonumber\\
 e_C&=\left\lfloor\frac{e_{\mathrm{in}}q_C}{q_{\mathrm{in}}}\right\rfloor,\nonumber\\
 e_T&=e_{\mathrm{in}}-e_C+\delta_T
 \label{eq:transfer}
\end{align}
남기는 몫을 질량에 비례해 정수로 계산하고 나머지는 보내는 몫에
준다. $\delta_T\geq0$는 운송 탄소 증분이다. 전량 이전에도
값이 모두 0인 잔여 Note를 만든다.

Voucher는 수신자의 인수 또는 송신자의 회수로 해결된다.
각자는 자기 소유권을 증명하고 같은 $rvnf$를 소비해
문서·역할·회계 값이 같은 자기 Note를 받는다. 먼저 수락된
경로만 성공한다.

회수는 기한이 지난 뒤가 아니라 \emph{기한까지} 허용한다.
공개 기한 $D$는 이전 시점보다 미래여야 하고, 회수 실행 시
\begin{equation}
 \mathsf{CurrentEpoch}()\leq D
 \label{eq:recall}
\end{equation}
를 만족해야 한다. 에포크는 블록 높이로 정한 시간 구간이며,
증명 생성 시각은 기준이 아니다. 인수에는 기한 제한이 없고,
두 경로 모두 실행 시점의 미소비·활성 검사를 거친다.

\subsection{생산 회계 규칙}

병합(Merge)은 소유자·문서 해시가 같은 일반용 Note 둘을 소비해
값을 합친 Note를 만든다. 출력의 소유자·문서·역할은 유지된다.
\begin{equation}
 (q,a,e)_{\mathrm{out}}=(q_1+q_2,a_1+a_2,e_1+e_2)
 \label{eq:merge}
\end{equation}
같은 입력을 두 번 넣는 것은 공통 중복 검사에서 거부한다.

분할(Split)은 질량이 양수인 일반용 Note를 같은 소유자·문서·
역할의 Note 둘로 나눈다.
$q_{\mathrm{in}}=q_1+q_2$일 때 배분은 다음과 같다.
\begin{align}
 a_2&=\left\lfloor\frac{a_{\mathrm{in}}q_2}{q_{\mathrm{in}}}\right\rfloor,
 &a_1&=a_{\mathrm{in}}-a_2,\nonumber\\
 e_2&=\left\lfloor\frac{e_{\mathrm{in}}q_2}{q_{\mathrm{in}}}\right\rfloor,
 &e_1&=e_{\mathrm{in}}-e_2
 \label{eq:split}
\end{align}
$\lfloor\cdot\rfloor$는 소수 부분을 버린다. 둘째 몫을 먼저 계산하고
나머지를 첫째 몫에 주므로 정수 잔여량이 보존된다.
질량 0 출력은 크레딧과 탄소도 0이어야 한다.

가공(Process)은 등록 정책에 따라 서로 다른 일반용 입력 $m$개를
출력 $n$개로 바꾼다. 입출력은 공정 수행자가 소유하고,
$m,n\geq1$의 지원 범위는 배포에서 정한다. 정책은 출력의 순서와
역할을 미리 고정하며 하나 이상을 일반용으로 지정한다.
이 절은 일반 설계의 규칙을 설명한다. 구현과 성능 측정에서 다룬
입출력 수는 \ref{sec:evaluation}절과 \ref{sec:implementation-limits}절에서 구분한다.

입력 질량은 일반용 출력, 폐기물 출력과 손실로 나뉜다.
$I_E,I_W$는 각각 일반용·폐기물 출력의 위치 집합이고,
$Q_{\mathrm{in}},Q_E,Q_W$는 입력·일반용 출력·폐기물 출력의 질량 합이다.
\begin{equation}
 \begin{gathered}
 Q_{\mathrm{in}}=Q_E+Q_W+q_{\mathrm{loss}},\\
 q_{\mathrm{loss}}\geq0,\qquad Q_E>0
 \end{gathered}
 \label{eq:mass}
\end{equation}
재활용 크레딧은 일반용 출력에만 배분한다.
\begin{equation}
 \sum_i a_i^{\mathrm{in}}=\sum_{j\in I_E}a_j^{\mathrm{out}}
 \label{eq:credit}
\end{equation}
각 출력의 $0\leq a_j\leq q_j$ 범위 안에서 배분하며 폐기물의 크레딧은 0이다.

탄소는 입력의 합에 공정 증분 $\delta_P\geq0$을 더하고,
일반용 출력의 질량에 비례해 나눈다. 정수 배분의 잔여량은
질량이 양수인 첫 일반용 출력 $j_0$에 준다. $Q_E>0$이므로
그런 출력이 적어도 하나 있다.
\begin{align}
 E&=\sum_i e_i^{\mathrm{in}}+\delta_P,\quad\delta_P\geq0,\nonumber\\
 e_j^{\mathrm{out}}&=\left\lfloor\frac{Eq_j^{\mathrm{out}}}{Q_E}\right\rfloor
     \quad(j\in I_E,\ j\ne j_0),\nonumber\\
 e_{j_0}^{\mathrm{out}}&=E-\sum_{j\in I_E\setminus\{j_0\}}e_j^{\mathrm{out}}
 \label{eq:process-carbon}
\end{align}
폐기물과 질량 0인 일반용 출력의 탄소는 0이다.
질량 0 출력의 역할과 위치는 유지한다. 폐기물의 실제 배출량이 0이라는
뜻은 아니다.
이 배분의 한계는 \ref{sec:limits}절에서 살펴본다.

출력 문서는 입력과 달라질 수 있다. 외부에서 준 질량 손실과
탄소 증분, 재료의 호환성이나 실제 공정 수행은 이 식만으로
확인할 수 없다. 정수 범위와 나눗셈 검사는 부록~\ref{app:formats}에 둔다.

\subsection{DPP 발급·검증}
\label{sec:dpp}

근거의 사용을 마치는 방법은 종료와 발급이다.
종료(Exit)는 일반용 또는 폐기물 증표를 소비하고 출력은 만들지 않는다.
실제 입력은 감사 기록에 남기며, 폐기물의 유일한 소비 경로다.

발급(Issue)은 질량이 양수인 일반용 증표를 소비하여 Claim을 만든다.
정책이 재활용 크레딧 비율의 하한을 $\tau_R=u_R/v_R$,
탄소 집약도의 상한을 $\tau_C=u_C/v_C$로 정하면
\begin{equation}
 v_Ra\geq u_Rq,\qquad v_Ce\leq u_Cq
 \label{eq:issue}
\end{equation}
를 검사한다. 분모는 양수이고 단위·배율과 정수 범위가 맞아야 한다.
이 검사는 양변에 분모를 곱해 나눗셈 없이 수행한다.
문서 해시 $d$·정책 $p$·공개 논스 $c$는 식~\eqref{eq:claim}에 연결된다.
입력 소비와 Claim·생성·감사 기록은 함께 갱신되므로,
같은 증표를 다시 발급하거나 발급 뒤 종료할 수 없다.
